\documentclass[11pt]{article}
\usepackage[T1]{fontenc}
\usepackage{amsmath,amssymb}
\usepackage[margin=30mm]{geometry}
\title{Minimal replacements for EPPA--III Claims 3.9 and 3.11}
\date{9 October 2026}
\begin{document}
\maketitle

\noindent
\textbf{Editorial rule.} These are suggested local substitutions in the
circulation manuscript; the manuscript source remains unchanged.

\subsection*{Claim 3.9 --- replace the incorrect neighbourhood sentence and binomial estimate}

Fix \(w\in H\setminus G\) and consider its profile
\(P(w)=(N_H(w)\cap C_i)_{i=1}^s\).
Let \(r\geq 1\) be the number of its coordinates of size \(k\).
As every automorphism of \(G\) extends to \(H\), the
\(\operatorname{Sym}(t)\wr\operatorname{Sym}(s)\)-orbit of \(P(w)\)
is realised by vertices of \(H\setminus G\). It contains at least
\(\binom sr\binom tk^r\geq s\binom tk\) profiles.
Since \(|H\setminus G|=st\), this gives \(\binom tk\leq t\),
and hence \(k\in\{1,t-1\}\).

\subsection*{Claim 3.11 --- replace only the final computational-check paragraph}

Since \(H\) is regular and both off-diagonal entries are matchings,
the two diagonal entries have the same bipartite degree.
If this degree is \(0\) or \(2\), interchanging \(D_1\) and \(D_2\)
gives case~(1). Otherwise all four entries are matchings.
After reordering the vertices within the four cliques we may assume
\[
 A=\begin{pmatrix}M&M\\ M&Z\end{pmatrix},
 \qquad Z\in\{M,\overline M\}.
\]
For \(Z=M\), one vertex of \(C_1\) has two common neighbours
with one vertex of \(C_2\), and no common neighbour with
the other. This contradicts EPPA, because exchanging the two
vertices of \(C_2\) is an automorphism of \(G\).
Thus \(Z=\overline M\), and \(H\) is the Wagner graph.

\subsection*{Verification and scope}

The last identification is literal: writing \(c_{i,a}\in C_i\) and
\(d_{j,a}\in D_j\), the Wagner graph has the eight-cycle
\[
 c_{1,0},\,d_{1,0},\,c_{2,0},\,d_{2,1},\,
 c_{1,1},\,d_{1,1},\,c_{2,1},\,d_{2,0},\,c_{1,0},
\]
and its four remaining edges join opposite vertices of this
cycle. For \(Z=M\), the two cross-nonedges from \(c_{1,0}\) to
\(c_{2,0}\) and \(c_{2,1}\) have, respectively, the common
neighbour sets \(\{d_{1,0},d_{2,0}\}\) and \(\varnothing\).

The first replacement uses the two-sided regularity of
Claim~3.8; the second uses the structural reduction of Claim~3.10
and the previously established \(s=t=2\). Neither of those
dependencies is repaired here.

\paragraph{Minor correction to the stated exceptional case.}
The matrix with four matching entries has degree \(3\), whereas
the complement of the Wagner graph has degree \(4\). Consequently
the phrase \emph{or its complement} is not justified by this
matrix description. The complemented result is obtained by
complementing both \(H\) and the embedded \(G\), rather than
by relabelling the same displayed matrix. Likewise,
complementing all four matching signs preserves their parity
and gives an isomorphic Wagner graph, \emph{not} the graph
complement.

\end{document}
